\section{Opacities}
\label{sec:opacities}

This module computes monochromatic and Rosseland-mean radiative opacities
for fully ionised hydrogen, covering both the non-magnetic ($B=0$) limit
and the strongly magnetised regime where the field is quantising
($\beta_e \equiv \hbar\omega_{\mathrm{ce}}/(k_B T) \gg 1$).
The magnetic treatment follows Potekhin \& Chabrier (2003, hereafter
PC03)\footnote{A.~Y.~Potekhin \& G.~Chabrier, \emph{ApJ} \textbf{585},
955 (2003).}, cross-referenced by equation number below.
The non-magnetic limit uses the Gaunt-factor--based formulae of Haakonsen
et~al.\ (2012)\footnote{C.~B.~Haakonsen et~al., \emph{ApJ} \textbf{749},
52 (2012).}.
Implementation resides in six source files under \texttt{src/opacity/}.


%======================================================================
\subsection{Non-magnetic opacities}
\label{sec:opacity:nonmag}
%======================================================================

\subsubsection{Free-free absorption}
\label{sec:opacity:nonmag:ff}

\noindent
The non-magnetic free-free absorption opacity per unit mass for fully
ionised hydrogen is
(\texttt{hydrogen\_ff.jl}, function \texttt{kappa\_ff}; Haakonsen et~al.\
Eq.~12):
%
\begin{equation}
\label{eq:kappa_ff_nonmag}
\kappa_\nu^{\mathrm{ff}}
  = \frac{4\,e^6}{3\,m_e\,h\,c}
    \sqrt{\frac{2\pi}{3\,k_B\,m_e}}\;
    \frac{\rho}{(m_p+m_e)^2}\;
    \nu^{-3}\;
    \bigl(1-e^{-u}\bigr)\;\tilde{g}_{\mathrm{ff}}(u,\gamma^2),
\end{equation}
%
where $u = h\nu/(k_B T)$ and $\gamma^2 = \mathrm{Ry}/(k_B T)$ with
Ry${}= 13.6057$~eV.  The thermally-averaged Gaunt factor
$\tilde{g}_{\mathrm{ff}}$ is obtained by bilinear interpolation on the
Sutherland (1998) table (\texttt{gaunt\_ff.jl}, struct
\texttt{GauntTable}).  Boundary values are clamped to the table edges,
reproducing the McPHAC convention.

\subsubsection{Thomson scattering}
\label{sec:opacity:nonmag:thomson}

\noindent
The reduced Thomson scattering opacity for fully ionised hydrogen
(Haakonsen Eq.~11 with $f_{\mathrm{ion}}=1$) is
%
\begin{equation}
\label{eq:sigma_thomson_nonmag}
\sigma_{\mathrm{T},t} = \frac{\sigma_T}{m_p+m_e},
\end{equation}
%
where $\sigma_T = 6.652\times10^{-25}$~cm$^2$ is the non-magnetic Thomson
cross-section.

\subsubsection{Total opacity and scattering albedo}
\label{sec:opacity:nonmag:total}

\noindent
The total extinction opacity is
%
\begin{equation}
\label{eq:k_total_nonmag}
k_\nu = \kappa_\nu^{\mathrm{ff}} + \sigma_{\mathrm{T},t},
\end{equation}
%
and the scattering albedo (Haakonsen Eq.~6) is
$\varrho_\nu = \sigma_{\mathrm{T},t}/k_\nu$.

\subsubsection{Non-magnetic Rosseland mean}
\label{sec:opacity:nonmag:ross}

\noindent
The Rosseland mean opacity (\texttt{hydrogen\_ff.jl}, function
\texttt{rosseland\_mean}) is computed by numerical quadrature on a
discrete frequency grid:
%
\begin{equation}
\label{eq:ross_nonmag}
\frac{1}{\kappa_R}
  = \frac{\displaystyle\int_0^\infty
      \frac{1}{k_\nu}\,\frac{\partial B_\nu}{\partial T}\,\mathrm{d}\nu}
    {\displaystyle\int_0^\infty
      \frac{\partial B_\nu}{\partial T}\,\mathrm{d}\nu},
\end{equation}
%
where the Planck derivative is
%
\begin{equation}
\label{eq:dBdT}
\frac{\partial B_\nu}{\partial T}
  = \frac{2h^2\nu^4}{c^2\,k_B\,T^2}\;\frac{e^u}{(e^u - 1)^2}.
\end{equation}
%
The integrals are evaluated by the midpoint rule on a log-spaced grid.
A fallback to $\sigma_{\mathrm{T},t}$ is returned when the denominator
vanishes.


%======================================================================
\subsection{Magnetic opacities: basic polarisation cross-sections}
\label{sec:opacity:mag:basic}
%======================================================================

In a magnetised plasma, radiative transfer decomposes into three basic
cyclic polarisations $\alpha\in\{-1,0,+1\}$ (left-circular, longitudinal,
right-circular with respect to~$\boldsymbol{B}$).  The code defines the
electron and proton cyclotron frequencies as
%
\begin{equation}
\label{eq:cyclotron}
\omega_{\mathrm{ce}} = \frac{eB}{m_e c}, \qquad
\omega_{\mathrm{cp}} = \frac{eB}{m_p c},
\end{equation}
%
and the quantisation parameter (PC03 Eq.~1)
%
\begin{equation}
\label{eq:beta_e}
\beta_e = \frac{\hbar\omega_{\mathrm{ce}}}{k_B T}.
\end{equation}

\subsubsection{Magnetic electron scattering}
\label{sec:opacity:mag:scat_e}

\noindent
The magnetic Thomson cross-section for polarisation~$\alpha$ is
(PC03 Eq.~33; \texttt{magnetic\_ff.jl}, function
\texttt{sigma\_scat\_alpha}):
%
\begin{equation}
\label{eq:sigma_scat_e}
\sigma_\alpha^{s,e}
  = \frac{\omega^2}{(\omega + \alpha\,\omega_{\mathrm{ce}})^2
         + \nu_{e,\alpha}^{s\,2}}\;\sigma_T,
\end{equation}
%
where the natural (radiative) damping frequency is (PC03 Eq.~39)
%
\begin{equation}
\label{eq:nu_e_s}
\nu_e^s = \frac{2}{3}\,\frac{e^2}{m_e\,c^3}\,\omega^2.
\end{equation}
%
In the non-magnetic limit ($B \to 0$) every polarisation receives
$\sigma_T/3$.

\subsubsection{Electron free-free absorption}
\label{sec:opacity:mag:ff}

\noindent
The magnetic free-free cross-section per polarisation is
(PC03 Eq.~51; \texttt{magnetic\_ff.jl}, function
\texttt{sigma\_ff\_alpha}):
%
\begin{equation}
\label{eq:sigma_ff}
\sigma_\alpha^{\mathrm{ff}}
  \approx \frac{\omega^2}
    {(\omega+\alpha\,\omega_{\mathrm{ce}})^2
     (\omega-\alpha\,\omega_{\mathrm{cp}})^2
     + \omega^2\,\tilde\nu_\alpha^2}\;
  \frac{4\pi e^2\,\nu_\alpha^{\mathrm{ff}}}{m_e\,c},
\end{equation}
%
where $\tilde\nu_\alpha$ is the combined damping frequency (PC03 Eq.~52):
%
\begin{equation}
\label{eq:nu_tilde}
\tilde\nu_\alpha
  = \nu_\alpha^{\mathrm{ff}}
  + \Bigl(1 + \alpha\,\frac{\omega_{\mathrm{ce}}}{\omega}\Bigr)\nu_p
  + \Bigl(1 - \alpha\,\frac{\omega_{\mathrm{cp}}}{\omega}\Bigr)\nu_e,
\end{equation}
%
and the effective free-free collision frequency is (PC03 Eq.~41):
%
\begin{equation}
\label{eq:nu_ff}
\nu_\alpha^{\mathrm{ff}}
  = \frac{4}{3}\sqrt{\frac{2\pi}{m_e\,k_B T}}\;
    \frac{n_e\,e^4}{\hbar\omega}\;
    (1-e^{-u})\;\Lambda_\alpha^{\mathrm{ff}},
  \qquad u \equiv \frac{\hbar\omega}{k_B T}.
\end{equation}
%
The factor $(1-e^{-u})$ accounts for induced (stimulated) radiation.

\paragraph{Note on damping terms.}
In Eq.~\eqref{eq:nu_tilde} the electron damping is
$\nu_e = \nu_e^s + \nu^{\mathrm{pp}}$
and the proton damping is
$\nu_p = \nu_p^s + \nu^{\mathrm{pp}}$,
where $\nu_p^s = (2/3)(e^2/(m_p c^3))\omega^2$ is the proton radiative
width (analogous to Eq.~\eqref{eq:nu_e_s}).
The code sets $\nu_e$ and $\nu_p$ from these sums before constructing
$\tilde\nu_\alpha$.  This is consistent with the structure of PC03
Eqs.~38--39 and~52, though PC03 do not write out an explicit
``$\nu_p$'' symbol.

\subsubsection{Proton-proton collisions}
\label{sec:opacity:mag:pp}

\noindent
Absorption due to proton-proton collisions is given by
(PC03 Eq.~46; \texttt{magnetic\_ff.jl}, function
\texttt{sigma\_pp\_alpha}):
%
\begin{equation}
\label{eq:sigma_pp}
\sigma_\alpha^{\mathrm{pp}}
  = \frac{1}{(\omega - \alpha\,\omega_{\mathrm{cp}})^2
         + \nu_{p,\alpha}^2}\;
    \frac{4\pi e^2\,\nu^{\mathrm{pp}}}{m_p\,c},
\end{equation}
%
where the proton damping frequency for this channel is
(PC03 Eq.~53):
%
\begin{equation}
\label{eq:nu_p_alpha}
\nu_{p,\alpha}
  = \frac{m_e}{m_p}\,\nu_\alpha^{\mathrm{ff}}
  + \nu_p^s + \nu^{\mathrm{pp}},
\end{equation}
%
and $\nu^{\mathrm{pp}}$ is the effective proton-proton collision frequency
(PC03 Eq.~47; \texttt{magnetic\_ff.jl}, function \texttt{nu\_pp}):
%
\begin{equation}
\label{eq:nu_pp}
\nu^{\mathrm{pp}}
  = \frac{256}{3}\sqrt{\frac{\pi}{m_p\,k_B T}}\;
    \frac{n_p\,e^4}{\hbar\omega}\;
    \frac{k_B T}{m_p c^2}\;
    (1-e^{-u})\;\Lambda_{\mathrm{pp}}.
\end{equation}

\subsubsection{Total polarisation cross-section}
\label{sec:opacity:mag:total_alpha}

\noindent
For fully ionised hydrogen ($x_H = 0$), the total cross-section per
polarisation is (\texttt{magnetic\_ff.jl}, function
\texttt{sigma\_total\_alpha}):
%
\begin{equation}
\label{eq:sigma_total_alpha}
\sigma_\alpha^{\mathrm{tot}}
  = \sigma_\alpha^{\mathrm{ff}}
  + \sigma_\alpha^{\mathrm{pp}}
  + \sigma_\alpha^{s,e}.
\end{equation}
%
This omits proton scattering $\sigma_\alpha^{s,p}$ (PC03 Eq.~34), which
scales as $(m_e/m_p)^2\approx 3\times 10^{-7}$ relative to
$\sigma_\alpha^{s,e}$ and is negligible away from the proton cyclotron
resonance.


%======================================================================
\subsection{Coulomb logarithms}
\label{sec:opacity:coulomb}
%======================================================================

\subsubsection{Classical (non-magnetic) Coulomb logarithm}
\label{sec:opacity:coulomb:classical}

\noindent
In the $B=0$ limit, the thermally averaged Coulomb logarithm is
(PC03 Eq.~43; \texttt{coulomb\_magnetic.jl}, function
\texttt{coulomb\_log\_classical\_safe}):
%
\begin{equation}
\label{eq:lambda_cl}
\Lambda_{\mathrm{cl}}^{\mathrm{ff}}
  = e^{u/2}\,K_0(u/2),
\end{equation}
%
where $K_0$ is the modified Bessel function of the second kind.  For
$u/2 > 300$ the code uses the asymptotic form
$e^x K_0(x) \sim \sqrt{\pi/(2x)}$ to avoid overflow.

\subsubsection{Magnetic Coulomb logarithm}
\label{sec:opacity:coulomb:magnetic}

\noindent
When the field is quantising ($\beta_e > 0.1$), the Born-approximation
Coulomb logarithm replaces $\Lambda_{\mathrm{cl}}^{\mathrm{ff}}$
(PC03 Eq.~44a; \texttt{coulomb\_magnetic.jl}, function
\texttt{coulomb\_log\_magnetic}):
%
\begin{equation}
\label{eq:lambda_mag}
\Lambda_\alpha^{\mathrm{ff}}
  = \frac{3}{4}\,e^{u/2}
    \sum_{n=-\infty}^{\infty}
    \int_0^\infty Q_n^\alpha(\beta_e,u,y)\;\mathrm{d}y.
\end{equation}
%
The integrand $Q_n^\alpha$ is given by (PC03 Eqs.~44b--e):
%
\begin{align}
Q_n^\alpha &= \frac{y}{\zeta}\;
  \bigl[(y+\theta+\zeta)\sinh(\beta_e/2)\bigr]^{-|n|}
  A_n^\alpha,
  \label{eq:Q_integrand}\\[4pt]
%
A_n^0 &= \frac{x_n\,K_1(x_n)}{y + \beta_e/4},
  \label{eq:A0}\\[4pt]
%
A_n^{\pm1} &= \frac{y + \theta + |n|\,\zeta}{\zeta^2}\;K_0(x_n),
  \label{eq:Apm1}
\end{align}
%
with auxiliary quantities (PC03 Eqs.~44d--e):
%
\begin{equation}
\label{eq:aux_mag_coulomb}
\zeta = \sqrt{1+2\theta y+y^2},\quad
\theta = \frac{1+e^{-\beta_e}}{1-e^{-\beta_e}},\quad
x_n = \frac{|u - n\beta_e|}{\sqrt{0.25 + y/\beta_e}}.
\end{equation}
%
The Landau-level sum is truncated at
$|n| \le N_{\max} = \min\!\bigl(\max(2,\lceil 3/\beta_e + 2\rceil),
50\bigr)$;
for strongly quantising fields ($\beta_e \gg 1$) the $n=0$ term
dominates.  Each integral is evaluated by adaptive quadrature (QuadGK)
over $y\in[0,20]$ with a relative tolerance of $10^{-4}$.

\subsubsection{Proton-proton Coulomb logarithm}
\label{sec:opacity:coulomb:pp}

\noindent
The proton-proton Coulomb logarithm is the analytic fit
(PC03 Eq.~48; \texttt{magnetic\_ff.jl}, function
\texttt{coulomb\_log\_proton}):
%
\begin{equation}
\label{eq:lambda_pp}
\Lambda_{\mathrm{pp}}
  \approx 0.6\,\ln\!\bigl(22\,u^{-1} + 9\,u^{-0.3}\bigr)
  + 0.4\,\sqrt{\pi u},
\end{equation}
%
which is stated to be accurate to ${\sim}1.5\%$ over the full range of~$u$.


%======================================================================
\subsection{Dielectric tensor and polarisation weights}
\label{sec:opacity:dielectric}
%======================================================================

\noindent
The mode-resolved opacity requires the polarisation weights
$|e_{j,\alpha}|^2$ for the two normal modes ($j=1$: extraordinary,
$j=2$: ordinary).  These are computed from the cold-plasma dielectric
tensor (\texttt{dielectric\_tensor.jl}) via the Stix parameters
(Ginzburg 1970; Shafranov 1967):
%
\begin{equation}
\label{eq:stix}
S = 1 - \frac{\omega_{pe}^2}{\omega^2-\omega_{\mathrm{ce}}^2}
      - \frac{\omega_{pi}^2}{\omega^2-\omega_{\mathrm{cp}}^2},\quad
D = \frac{\omega_{\mathrm{ce}}\,\omega_{pe}^2}
         {\omega(\omega^2-\omega_{\mathrm{ce}}^2)}
  - \frac{\omega_{\mathrm{cp}}\,\omega_{pi}^2}
         {\omega(\omega^2-\omega_{\mathrm{cp}}^2)},\quad
P = 1 - \frac{\omega_{pe}^2 + \omega_{pi}^2}{\omega^2},
\end{equation}
%
where $\omega_{pe}^2 = 4\pi n_e e^2/m_e$ and
$\omega_{pi}^2 = 4\pi n_e e^2/m_p$.
The mode parameter $q$ (PC03 Eq.~25) is
%
\begin{equation}
\label{eq:q_param}
q = \frac{(P-S)\sin^2\theta_B}{2\,D\cos\theta_B},
\end{equation}
%
where $\theta_B$ is the angle between the photon propagation direction
and~$\boldsymbol{B}$.  The transverse polarisation ratios follow from the
quadratic $K^2 + 2qK - 1 = 0$ (Shafranov 1967):
%
\begin{equation}
\label{eq:K_modes}
K_1 = -q + \sqrt{1+q^2}\quad\text{(extraordinary)},\qquad
K_2 = -q - \sqrt{1+q^2}\quad\text{(ordinary)}.
\end{equation}
%
Note that $K_1 K_2 = -1$.  The longitudinal component is
%
\begin{equation}
\label{eq:K_z}
K_{z,j} = -\frac{K_j\,(S - K_j D)\,\sin\theta_B\cos\theta_B}
               {P - (S - K_j D)\sin^2\theta_B},
\end{equation}
%
from which the cyclic weights are extracted:
%
\begin{equation}
\label{eq:weights}
|e_{j,+1}|^2 = \frac{(K_j+1)^2}{2\mathcal{N}_j},\qquad
|e_{j,-1}|^2 = \frac{(K_j-1)^2}{2\mathcal{N}_j},\qquad
|e_{j,0}|^2  = \frac{K_{z,j}^2}{\mathcal{N}_j},
\end{equation}
%
with normalisation $\mathcal{N}_j = 1 + K_j^2 + K_{z,j}^2$.
The weights satisfy $\sum_\alpha |e_{j,\alpha}|^2 = 1$ after a safety
renormalisation clamp.

\paragraph{Limiting cases.}
At $\theta_B = 0$ the modes are purely circularly polarised:
mode~1 has all weight on $\alpha=+1$, mode~2 on $\alpha=-1$.
At $\theta_B = \pi/2$ mode~1 is purely longitudinal ($\alpha=0$) and
mode~2 is an equal mix of $\alpha=\pm1$.
These limits are hard-coded as special cases to avoid division by zero
in Eq.~\eqref{eq:q_param}.

\paragraph{Design choice.}
An earlier version of this code used the approximate relation
$q(\omega,\theta_B) \approx \bar q(\omega)\sin^2\theta_B/(2\cos\theta_B)$
(PC03 Eq.~26), which is inaccurate near cyclotron resonances.
The current implementation uses the full Stix-parameter formulation
of Eq.~\eqref{eq:q_param}, which is numerically stable at all densities
because $q$ is determined primarily by $\omega_{\mathrm{ce}}/\omega$.


%======================================================================
\subsection{Normal-mode opacities and the absorption/scattering split}
\label{sec:opacity:modes}
%======================================================================

\noindent
Given the polarisation weights, the opacity for normal mode~$j$ at
angle~$\theta_B$ is (PC03 Eq.~27; \texttt{magnetic\_modes.jl}, function
\texttt{mode\_opacity}):
%
\begin{equation}
\label{eq:kappa_mode}
\kappa_j(\theta_B)
  = \frac{1}{m_H} \sum_{\alpha=-1}^{+1}
    |e_{j,\alpha}(\theta_B)|^2\;\sigma_\alpha^{\mathrm{tot}},
\end{equation}
%
where $m_H = m_p + m_e$ and $\sigma_\alpha^{\mathrm{tot}}$ is given by
Eq.~\eqref{eq:sigma_total_alpha}.

The code provides a clean \emph{absorption/scattering split} for use by
the Feautrier transport solver, which requires them separately:
%
\begin{align}
\kappa_j^{\mathrm{abs}}(\theta_B)
  &= \frac{1}{m_H}\sum_\alpha |e_{j,\alpha}|^2\,
     \bigl(\sigma_\alpha^{\mathrm{ff}}
     + \sigma_\alpha^{\mathrm{pp}}\bigr),
  \label{eq:mode_absorption}\\[4pt]
\kappa_j^{\mathrm{scat}}(\theta_B)
  &= \frac{1}{m_H}\sum_\alpha |e_{j,\alpha}|^2\,
     \sigma_\alpha^{s,e}.
  \label{eq:mode_scattering}
\end{align}
%
These are implemented as the functions \texttt{mode\_absorption} and
\texttt{mode\_scattering} in \texttt{magnetic\_modes.jl}.
By construction,
$\kappa_j = \kappa_j^{\mathrm{abs}} + \kappa_j^{\mathrm{scat}}$.

\paragraph{Effective (unpolarised) opacity.}
For unpolarised radiation at a single angle, the effective opacity is the
harmonic mean of the two normal modes:
%
\begin{equation}
\label{eq:kappa_eff}
\kappa_{\mathrm{eff}} = \frac{2}{1/\kappa_1 + 1/\kappa_2}.
\end{equation}


%======================================================================
\subsection{Angle-averaged opacities and Rosseland means}
\label{sec:opacity:rosseland_mag}
%======================================================================

\noindent
The diffusion approximation for transport along ($\parallel$) and across
($\perp$) the magnetic field introduces angle-averaged opacities for each
mode separately (PC03 Eq.~30; \texttt{magnetic\_modes.jl}, function
\texttt{kappa\_parallel\_perp\_mono}):
%
\begin{equation}
\label{eq:angle_avg}
\frac{1}{\kappa_j^\parallel}
  = \frac{3}{4}\int_0^\pi
    \frac{\cos^2\theta_B}{\kappa_j(\theta_B)}\,
    \sin\theta_B\;\mathrm{d}\theta_B,
\qquad
\frac{1}{\kappa_j^\perp}
  = \frac{3}{2}\int_0^\pi
    \frac{\sin^3\theta_B}{\kappa_j(\theta_B)}\;
    \mathrm{d}\theta_B.
\end{equation}
%
These integrals are evaluated by the midpoint rule with $N_\theta = 40$
uniform panels.  After obtaining $\kappa_j^\parallel$ and $\kappa_j^\perp$
for each mode $j\in\{1,2\}$, the code forms the harmonic mean over modes:
%
\begin{equation}
\label{eq:harmonic_modes}
\kappa_{\mathrm{eff}}^\parallel
  = \frac{2}{1/\kappa_1^\parallel + 1/\kappa_2^\parallel},
\qquad
\kappa_{\mathrm{eff}}^\perp
  = \frac{2}{1/\kappa_1^\perp + 1/\kappa_2^\perp}.
\end{equation}

\paragraph{Importance of mode-separate averaging.}
A previous version of the code incorrectly combined the two normal modes
\emph{before} angle-averaging, which overpredicts the opacity when the two
modes differ by orders of magnitude (as they typically do near the
cyclotron resonance).  The current implementation performs the
angle-average on each mode independently, then combines them, following
the prescription of PC03 Eq.~30.

\subsubsection{Magnetic Rosseland means}

\noindent
The Rosseland mean opacities for transport along and across the field are
(\texttt{magnetic\_modes.jl}, function \texttt{rosseland\_magnetic}):
%
\begin{equation}
\label{eq:ross_mag}
\frac{1}{K_\parallel}
  = \frac{\displaystyle\int_0^\infty
      \frac{1}{\kappa_{\mathrm{eff}}^\parallel(\nu)}\,
      \frac{\partial B_\nu}{\partial T}\,\mathrm{d}\nu}
    {\displaystyle\int_0^\infty
      \frac{\partial B_\nu}{\partial T}\,\mathrm{d}\nu},
\qquad
\frac{1}{K_\perp}
  = \frac{\displaystyle\int_0^\infty
      \frac{1}{\kappa_{\mathrm{eff}}^\perp(\nu)}\,
      \frac{\partial B_\nu}{\partial T}\,\mathrm{d}\nu}
    {\displaystyle\int_0^\infty
      \frac{\partial B_\nu}{\partial T}\,\mathrm{d}\nu}.
\end{equation}
%
These correspond to columns~13 and~14 of the Potekhin opacity tables
(PC03, Table~1).

The frequency grid is constructed by \texttt{make\_rosseland\_frequency\_grid}:
$N_{\mathrm{base}}=200$ logarithmically spaced points spanning
$[0.01,\,200]\times k_B T/h$, augmented with $N_{\mathrm{res}}=20$
additional points clustered within $\pm 0.3$~dex of each cyclotron
resonance ($\nu_{\mathrm{ce}}$ and~$\nu_{\mathrm{cp}}$) to resolve the
sharp features in the cross-sections near these frequencies.


%======================================================================
\subsection{Summary of equation cross-references}
\label{sec:opacity:xref}
%======================================================================

Table~\ref{tab:opacity_xref} maps every implemented equation to
its source in PC03 and to the corresponding Julia function.

\begin{table}[ht]
\centering
\caption{Cross-reference of opacity equations to Potekhin \& Chabrier
(2003) and source code.}
\label{tab:opacity_xref}
\small
\begin{tabular}{lllp{5.5cm}}
\hline
Report Eq. & PC03 Eq. & Source file / function & Description \\
\hline
\eqref{eq:beta_e}           & (1)     & \texttt{magnetic\_ff / beta\_e} &
  Quantisation parameter $\beta_e$ \\
\eqref{eq:q_param}          & (25)    & \texttt{dielectric\_tensor / polarization\_weights\_full} &
  Mode parameter $q$ \\
\eqref{eq:kappa_mode}        & (27)    & \texttt{magnetic\_modes / mode\_opacity} &
  Mode opacity $\kappa_j(\theta_B)$ \\
\eqref{eq:angle_avg}        & (30)    & \texttt{magnetic\_modes / kappa\_parallel\_perp\_mono} &
  Angle-averaged $\kappa_j^{\parallel,\perp}$ \\
\eqref{eq:sigma_scat_e}     & (33)    & \texttt{magnetic\_ff / sigma\_scat\_alpha} &
  Magnetic Thomson $\sigma_\alpha^{s,e}$ \\
\eqref{eq:sigma_ff}         & (37,51) & \texttt{magnetic\_ff / sigma\_ff\_alpha} &
  Free-free cross-section $\sigma_\alpha^{\mathrm{ff}}$ \\
\eqref{eq:nu_e_s}           & (39)    & \texttt{magnetic\_ff / sigma\_ff\_alpha} (inline) &
  Radiative damping $\nu_e^s$ \\
\eqref{eq:nu_ff}            & (41)    & \texttt{magnetic\_ff / nu\_ff\_alpha} &
  Free-free collision freq.\ $\nu_\alpha^{\mathrm{ff}}$ \\
\eqref{eq:lambda_cl}        & (43)    & \texttt{coulomb\_magnetic / coulomb\_log\_classical\_safe} &
  Classical Coulomb log $\Lambda_{\mathrm{cl}}^{\mathrm{ff}}$ \\
\eqref{eq:lambda_mag}       & (44a)   & \texttt{coulomb\_magnetic / coulomb\_log\_magnetic} &
  Magnetic Coulomb log $\Lambda_\alpha^{\mathrm{ff}}$ \\
\eqref{eq:Q_integrand}--\eqref{eq:aux_mag_coulomb}
                             & (44b--e)& \texttt{coulomb\_magnetic / \_Q\_n\_alpha} &
  Integrand components \\
\eqref{eq:sigma_pp}         & (46)    & \texttt{magnetic\_ff / sigma\_pp\_alpha} &
  Proton-proton $\sigma_\alpha^{\mathrm{pp}}$ \\
\eqref{eq:nu_pp}            & (47)    & \texttt{magnetic\_ff / nu\_pp} &
  Proton collision freq.\ $\nu^{\mathrm{pp}}$ \\
\eqref{eq:lambda_pp}        & (48)    & \texttt{magnetic\_ff / coulomb\_log\_proton} &
  Proton Coulomb log $\Lambda_{\mathrm{pp}}$ \\
\eqref{eq:nu_tilde}         & (52)    & \texttt{magnetic\_ff / sigma\_ff\_alpha} (inline) &
  Combined damping $\tilde\nu_\alpha$ \\
\eqref{eq:nu_p_alpha}       & (53)    & \texttt{magnetic\_ff / sigma\_pp\_alpha} (inline) &
  Proton damping $\nu_{p,\alpha}$ \\
\hline
\end{tabular}
\end{table}


%======================================================================
\subsection{Discrepancies and implementation notes}
\label{sec:opacity:discrepancies}
%======================================================================

\begin{enumerate}

\item \textbf{Proton scattering omitted.}
  PC03 Eq.~(34) gives the proton scattering cross-section
  $\sigma_\alpha^{s,p} = (m_e/m_p)^2\,
  \omega^2/[(\omega-\alpha\omega_{\mathrm{cp}})^2 + \nu_{p,\alpha}^2]
  \,\sigma_T$.
  This term is not included in \texttt{sigma\_total\_alpha}
  (Eq.~\ref{eq:sigma_total_alpha}).
  The omission is justified: away from the proton cyclotron resonance,
  $\sigma_\alpha^{s,p}/\sigma_\alpha^{s,e} \sim (m_e/m_p)^2
  \approx 3\times10^{-7}$.
  At $\omega \approx \omega_{\mathrm{cp}}$ the proton resonance
  is already captured in $\sigma_\alpha^{\mathrm{pp}}$
  (Eq.~\ref{eq:sigma_pp}), so the missing scattering contribution
  affects only the scattering-vs-absorption decomposition at that
  resonance, not the total opacity.

\item \textbf{Traditional vs.\ improved free-free formula.}
  PC03 note that the traditional factorisation (their Eq.~45,
  $\sigma_\alpha^{\mathrm{a}} = \sigma_\alpha^{\mathrm{ff}} /
  \sigma_\alpha^{s,e}$) is ``erroneous'' because it misses
  interference between terms with changing and fixed Landau quantum
  number~$n$.
  Our code implements the improved Eq.~(51) directly
  (Eq.~\ref{eq:sigma_ff}), which correctly incorporates the
  proton-cyclotron suppression factor $(\omega-\alpha\omega_{\mathrm{cp}})^2$
  in the denominator.

\item \textbf{$\nu_\alpha^{\mathrm{ff}}$ uses $\hbar\omega$, not $h\nu$.}
  The magnetic module defines $u = \hbar\omega/(k_B T)$ in
  Eq.~\eqref{eq:nu_ff}, whereas the non-magnetic module
  (\texttt{hydrogen\_ff.jl}) uses $u = h\nu/(k_B T)$.
  Since $\hbar\omega = h\nu$, these are identical; however, the
  intermediate code in \texttt{nu\_ff\_alpha} divides by $\hbar\omega$
  (not $h\nu$), which is consistent with the convention of PC03 Eq.~41
  where frequencies appear as angular frequencies~$\omega$.

\item \textbf{Coulomb logarithm switching threshold.}
  The code switches from the classical $\Lambda_{\mathrm{cl}}^{\mathrm{ff}}$
  to the magnetic $\Lambda_\alpha^{\mathrm{ff}}$ at $\beta_e = 0.1$.
  PC03 state that the non-magnetic Gaunt factor is ``applicable if
  $\beta_e < 1$'' (their text below Eq.~43).  The code's more conservative
  threshold of $\beta_e = 0.1$ ensures that the magnetic expression is
  used whenever the field produces even modest Landau quantisation, at
  the cost of evaluating the more expensive sum+integral for weakly
  quantising fields.

\item \textbf{Integral truncation in $\Lambda_\alpha^{\mathrm{ff}}$.}
  The semi-infinite integral $\int_0^\infty Q_n^\alpha\,\mathrm{d}y$ in
  Eq.~\eqref{eq:lambda_mag} is truncated at $y=20$ in the code.
  For $\beta_e \gg 1$ the integrand decays exponentially for $y \gg 1$
  (via the Bessel functions $K_0,\,K_1$ of argument~$x_n$), so this
  upper bound is adequate.  For marginally quantising fields
  ($\beta_e \sim 0.1$--$1$) the decay is slower but the code falls back
  to $\Lambda_{\mathrm{cl}}^{\mathrm{ff}}$ at $\beta_e < 0.1$, so the
  truncation error is bounded.

\item \textbf{Gaunt factor source.}
  The non-magnetic Gaunt factor table
  (\texttt{gaunt\_ff.jl}) interpolates the Sutherland (1998) numerical
  table rather than the Hummer (1988) Pad\'e fit used in PC03.
  For free-free emission by hydrogen at the temperatures of interest
  ($T \sim 10^{5.3}$--$10^{7}$~K), the two sources agree to better
  than 1\%, so this difference is inconsequential.

\end{enumerate}
